%% ma: L11-B15-0088
%% lop: 11
%% chuong: 5
%% bai: 15
%% dang: 11.15.7
%% loai: TLN
%% muc: TH
%% dap_an: "1839"
%% nguon: "tai-lieu-hoc-tap-mon-toan-lop-11-theo-chuong-trinh-giao-duc-pho-thong-2018.pdf (Lớp Toán Thầy Duy, Phạm Lê Duy, THPT Chu Văn An), Chương 3 Bài 1, Phần 3, Câu 10 (tr.343)"
%% kien_thuc: "Lớp 11 Bài 15 (giới hạn của dãy số)"
%% trang_thai: chinh-thuc
%% ghi_chu: "Đã duyệt 03/11/2026."
\begin{ex}
Số thập phân vô hạn tuần hoàn $0{,}271414\ldots$ viết dạng phân số có dạng $\dfrac mn$ với $m;n$ là các số tự nhiên và $\dfrac mn$ là phân số tối giản. Tính $n-3m$.
\shortans{1839}
\loigiai{\sloppy \[0{,}271414\ldots=0{,}27+0{,}001414\ldots=\dfrac{27}{100}+\dfrac{0{,}0014}{1-0{,}01}=\dfrac{27}{100}+\dfrac{14}{9900}=\dfrac{2687}{9900}.\] Vì $2687$ không chia hết cho $2,3,5,11$ nên phân số tối giản: $m=2687$, $n=9900$, và $n-3m=9900-8061=1839$.}
\end{ex}
